% Generated by paper/generate_data.py from the reviewed finite inputs.
\section{Toric and enumerative data}\label{app:data}

The source rays, numbered from zero, are the rows of the first matrix
below. The columns of $N$ give $H_1,\ldots,H_6$ in the divisor basis
$(D_2,D_5,D_6,D_7,D_8,D_9)$; $\det N=-1$. These coordinates and
the period conventions agree with \cite[Appendix A]{WuebbenHiggs}.
\[
 V=\begin{pmatrix}
1&0&0&0\\
0&1&0&0\\
-1&-1&0&0\\
0&0&1&0\\
0&0&0&1\\
-4&-2&-1&0\\
-4&-2&0&-1\\
3&1&1&1\\
2&1&1&1\\
-2&-1&0&0\\
-1&0&0&0
\end{pmatrix},\qquad N=\begin{pmatrix}
0&0&1&1&2&2\\
1&1&2&4&5&5\\
0&1&2&4&4&5\\
0&-1&-1&-3&-3&-3\\
0&-1&-1&-2&-2&-2\\
0&0&1&2&2&2
\end{pmatrix}.
\]
The maximal source cones are the following tuples of ray indices:\[\begin{gathered}
(0,1,3,6),\quad (0,1,3,7),\quad (0,1,4,5),\quad (0,1,4,7)\\
(0,1,5,6),\quad (0,2,4,5),\quad (0,2,4,7),\quad (0,2,5,6)\\
(0,2,6,7),\quad (0,3,6,7),\quad (1,3,6,10),\quad (1,3,7,8)\\
(1,3,8,10),\quad (1,4,5,10),\quad (1,4,7,8),\quad (1,4,8,10)\\
(1,5,6,10),\quad (2,4,5,9),\quad (2,4,7,9),\quad (2,5,6,9)\\
(2,6,7,9),\quad (3,6,7,9),\quad (3,6,9,10),\quad (3,7,8,9)\\
(3,8,9,10),\quad (4,5,9,10),\quad (4,7,8,9),\quad (4,8,9,10)\\
(5,6,9,10).
\end{gathered}\]
The fan is regular for heights
$(10,11,3,9,2,7,1,4,8,-5,-3)_i+(i+1)/1009$.
The charge matrix, with columns $0,\ldots,9$, is
\[Q=\begin{pmatrix}
0&0&-1&0&-1&1&0&1&0&0\\
0&0&0&1&1&0&0&0&-1&-1\\
0&1&0&0&0&0&0&2&-2&1\\
1&0&-1&0&0&0&0&0&0&1\\
0&0&1&1&0&0&-1&-1&0&0\\
0&0&0&-1&0&0&1&0&1&-1
\end{pmatrix}.\]

The symmetric cubic $\kappa_{abc}=H_aH_bH_c$ is specified by the
following entries with $a\le b\le c$. Indices here run from one to six.
\begin{center}\small
\begin{tabular}{rr@{\qquad}rr@{\qquad}rr}
\toprule
$abc$&$\kappa_{abc}$&$abc$&$\kappa_{abc}$&$abc$&$\kappa_{abc}$\\
\midrule
111&0&156&20&335&31\\
112&0&166&24&336&35\\
113&0&222&0&344&49\\
114&0&223&2&345&57\\
115&0&224&3&346&65\\
116&0&225&4&355&65\\
122&0&226&4&356&75\\
123&2&233&6&366&85\\
124&3&234&11&444&83\\
125&4&235&13&445&101\\
126&4&236&15&446&115\\
133&4&244&17&455&117\\
134&8&245&22&456&136\\
135&8&246&25&466&155\\
136&10&255&26&555&133\\
144&14&256&30&556&156\\
145&16&266&34&566&180\\
146&19&333&14&666&204\\
155&16&334&27&&\\
\bottomrule
\end{tabular}\end{center}

Table~\ref{tab:characters} lists sections of $2\mathcal D$ grouped
by their character. Parenthesized tuples denote cones. The 41 listed
cones are maximal among those with nonsingleton Newton face for the
pulled-back original equation. Their faces comprise all 99 cones whose
orbits may meet the hypersurface in the subdivided fan. The inequalities in Section~\ref{sec:geometry} can
be checked directly against its fifteen rays.
\begingroup\small
\begin{longtable}{p{0.25\textwidth}p{0.67\textwidth}}
\caption{Characters generating $2\mathcal D$ along the hypersurface.}
\label{tab:characters}\\
\toprule Character & Cones\\\midrule\endfirsthead
\toprule Character & Cones\\\midrule\endhead
$(0,0,0,10)$ & $(0,1,3)$, $(0,3,7)$, $(1,3,7)$, $(1,3,8)$, $(3,7,8)$\\[3pt]
$(0,0,10,0)$ & $(0,1,4)$, $(0,1,7)$, $(0,4,7)$, $(1,4,7)$, $(1,4,8)$, $(1,7,8)$, $(4,7,8)$\\[3pt]
$(0,0,14,14)$ & $(0,1,5)$, $(0,1,6)$, $(1,5,6)$\\[3pt]
$(0,4,6,2)$ & $(0,2,5)$, $(2,5,11)$\\[3pt]
$(0,4,6,6)$ & $(0,2,6)$, $(0,5,6)$, $(2,5,6)$\\[3pt]
$(0,4,2,4)$ & $(0,2,7)$, $(2,7,12)$\\[3pt]
$(0,0,0,14)$ & $(0,3,6)$\\[3pt]
$(0,2,10,0)$ & $(0,4,5)$, $(4,5,11)$\\[3pt]
$(2,0,0,6)$ & $(1,3,6)$, $(3,6,12)$, $(3,6,14)$, $(3,8,14)$\\[3pt]
$(0,0,14,0)$ & $(1,4,5)$\\[3pt]
$(0,4,2,6)$ & $(2,6,12)$\\[3pt]
$(0,4,4,2)$ & $(2,7,11)$\\[3pt]
$(0,2,0,8)$ & $(3,7,12)$\\[3pt]
$(2,0,6,0)$ & $(4,5,13)$, $(4,8,13)$\\[3pt]
$(0,2,8,0)$ & $(4,7,11)$\\[3pt]
$(3,0,2,2)$ & $(5,6,9)$, $(5,9,13)$, $(6,9,14)$, $(8,9,13)$, $(8,9,14)$\\[3pt]
\bottomrule\end{longtable}\endgroup

The smoothing rays are the rows of
\[
\begin{pmatrix}
-1&-3&-1&0\\-1&-3&0&-1\\-1&-1&0&0\\-1&0&0&0\\
0&1&0&0\\1&-1&0&0\\1&3&1&1
\end{pmatrix}.
\]
Their maximal cones are
$1456,1346,0456,0346,0145,0134,1236,0236,0123,01256$.
The last cone is nonsimplicial. The divisor charge matrix in the
basis $K_1,K_2,K_3$ is
\[
\begin{pmatrix}
1&1&-2&0&0&-1&1\\0&0&-1&2&0&1&0\\0&0&1&-1&1&0&0
\end{pmatrix}.
\]

\subsection*{Reproducibility}
The accompanying project retains the exact fan, monomial-section,
period and coefficient data. From its root, the source numerical
calculation is reproduced by
\begin{quote}\ttfamily\small
sage -python checks/x9\_source\_genus\_zero\_box.py\par
\quad --source-root ../global\_higgs\_branch\par
\quad --output /tmp/x9-source-genus-zero.json
\end{quote}
The source files used for the charge and cubic inputs are pinned by
SHA-256. The finite quotient has 431 positive monomials. All six
equations in \eqref{eq:mirror} have zero remainder; the computation
also checks 1440 Gamma contiguity identities. These arithmetic
checks are separate from the mirror-theorem hypotheses in
Section~\ref{sec:calculation}.

The scripts \texttt{checks/x9\_lift\_bound.py} and
\texttt{checks/x9\_contracted\_support.py}, run with Sage, reproduce
the finite lift and contracted-locus calculations.
The standard-library scripts \texttt{checks/x9\_plane\_geometry.py}
and \texttt{checks/x9\_universal\_line.py} reproduce the smoothing
fan calculation and the 37-row universal table. The stored exact
outputs are in \texttt{evidence/}. The symbolic two-by-two matrix
comparison for \eqref{eq:discrepancy} is
\texttt{checks/local\_series\_comparison.py}, using SymPy.
